\documentclass[a4paper,11pt]{article}
\usepackage[a4paper,margin=2cm]{geometry}
\usepackage[T1]{fontenc}
\usepackage[utf8]{inputenc}
\usepackage[ngerman,english]{babel}
\usepackage{lmodern,microtype,xcolor,parskip,enumitem}
\usepackage[most]{tcolorbox}
\usepackage{titlesec}
\usepackage[hidelinks]{hyperref}
\definecolor{HeaderBlueDark}{RGB}{70,110,170}
\definecolor{RowBlueLight}{RGB}{235,243,255}
\definecolor{RowBlueDark}{RGB}{220,233,250}
\definecolor{SoftGray}{RGB}{90,90,90}
\setlength{\parindent}{0pt}
\setlist[itemize]{leftmargin=1.4em,itemsep=0.35em,topsep=0.35em}
\titleformat{\section}[block]{\Large\bfseries\color{HeaderBlueDark}}{}{0pt}{}
\newtcolorbox{exbox}{enhanced,breakable,colback=RowBlueLight,colframe=RowBlueDark,boxrule=0.4pt,arc=1.5mm,left=1.2mm,right=1.2mm,top=1mm,bottom=1mm,title={Beispiele \textnormal{(Examples)}},fonttitle=\bfseries,colbacktitle=RowBlueDark,coltitle=black}
\NewDocumentEnvironment{grammarentry}{mm+b}{\begin{tcolorbox}[enhanced,breakable,colback=white,colframe=HeaderBlueDark,title={#1\hfill\normalfont\footnotesize #2},colbacktitle=HeaderBlueDark,coltitle=white,fonttitle=\bfseries]#3\end{tcolorbox}}{}
\newcommand{\GermanText}[1]{\textbf{Deutsch: }#1\par}
\newcommand{\EnglishText}[1]{{\small\color{SoftGray}\textbf{English: }#1\par}}
\newcommand{\ExampleItem}[2]{\item \textbf{#1}\\{\small\color{SoftGray}#2}}
\title{Die Modalpartikel \glqq etwa\grqq\\{\large\textnormal{The modal particle etwa}}}
\date{}
\begin{document}
\maketitle
\tableofcontents
\section{Bedeutung und Wirkung \textnormal{(Meaning and effect)}}
\begin{grammarentry}{etwa}{Modalpartikel / Abtönungspartikel}
\GermanText{Etwa bezeichnet in Fragen eine unerwartete, häufig unerwünschte oder beunruhigende Möglichkeit.}
\EnglishText{Etwa in questions signals an unexpected, often unwelcome or worrying possibility.}
\end{grammarentry}
\section{Grammatische Rolle und Satzarten \textnormal{(Grammar and sentence moods)}}
\begin{grammarentry}{Verwendung}{Satzart und Kontext}
\GermanText{Vor allem in Ja/Nein-Fragen, gelegentlich zusammen mit denn und nicht. Der Sprecher hatte oft das Gegenteil erwartet.}
\EnglishText{Mainly yes/no questions, sometimes combined with denn and nicht. The speaker often expected the opposite.}
\end{grammarentry}
\begin{grammarentry}{Form und Stellung}{Syntax}
\GermanText{Als Modalpartikel ist etwa unflektierbar, fordert keinen Kasus und steht normalerweise im Mittelfeld. Allein im Vorfeld eines normalen Aussagesatzes kann eine Modalpartikel nicht stehen. Die Betonung und der Kontext beeinflussen die Wirkung.}
\EnglishText{As a modal particle, etwa is invariable, does not govern case and normally appears in the middle field. It cannot occupy the prefield of a normal statement on its own. Stress and context affect its pragmatic force.}
\end{grammarentry}
\section{Typische Beispiele \textnormal{(Typical examples)}}
\begin{exbox}\begin{itemize}
\ExampleItem{Hast du etwa Angst?}{Don't tell me you're afraid?}
\ExampleItem{Willst du etwa aufgeben?}{You're not thinking of giving up, are you?}
\ExampleItem{Kommst du etwa nicht?}{Don't tell me you're not coming?}
\ExampleItem{Hast du etwa nicht gelernt?}{Don't tell me you didn't study?}
\ExampleItem{Hast du denn etwa alles vergessen?}{Have you seriously forgotten everything?}
\end{itemize}\end{exbox}
\section{Abgrenzung zu anderen Verwendungen \textnormal{(Contrast with other uses)}}
\begin{grammarentry}{Wortarten unterscheiden}{Nicht immer Modalpartikel}
\GermanText{Als Mengenangabe heißt etwa ungefähr: etwa zwanzig Euro. Diese Näherungsangabe ist keine Modalpartikel.}
\EnglishText{As a quantity expression, etwa means approximately: about twenty euros. This is not a modal-particle use.}
\end{grammarentry}
\begin{exbox}\begin{itemize}
\ExampleItem{Es kostet etwa zwanzig Euro.}{It costs about twenty euros.}
\ExampleItem{Kommst du nicht?}{Aren't you coming? (more neutral)}
\end{itemize}\end{exbox}
\section{etwa nicht \textnormal{(etwa with nicht)}}
\begin{grammarentry}{Wichtiger Unterschied}{Important distinction}
\GermanText{In Kommst du etwa nicht? ist etwa eine Modalpartikel und nicht gewöhnlich eine Negationspartikel. Die Kombination ist keine selbstständige Modalpartikel. Etwa drückt die überraschende Erwartung aus, nicht verneint die Aussage.}
\EnglishText{In Kommst du etwa nicht?, etwa is a modal particle and nicht is normally a negation particle. Together they are not a separate modal particle. Etwa adds an attitude of surprise; nicht negates.}
\end{grammarentry}
\begin{exbox}\begin{itemize}
\ExampleItem{Kommst du denn etwa nicht?}{So, are you seriously not coming?}
\ExampleItem{Hat sie etwa noch nicht angerufen?}{Don't tell me she still hasn't called?}
\end{itemize}\end{exbox}
\section{Lernhinweis \textnormal{(Learning tip)}}
\GermanText{Eine Modalpartikel verändert vor allem die Gesprächsabsicht. Versuchen Sie nicht, jedem Vorkommen eine feste englische Eins-zu-eins-Übersetzung zu geben. Achten Sie auf Satzart, Tonfall und Situation.}
\EnglishText{A modal particle mainly affects the speaker's intention. Do not try to give it one fixed English translation in every context. Pay attention to sentence mood, intonation and situation.}
\section{Quellen \textnormal{(Sources)}}
\begin{itemize}
\item IDS Mannheim: \url{https://grammis.ids-mannheim.de/terminologie/2}
\item Duden: \url{https://www.duden.de/sprachwissen/sprachratgeber/Modalpartikeln}
\end{itemize}
\end{document}
